% ---------------------------------------------------------------
% preamble/lang-th.tex - ฟอนต์ไทย การตัดคำ และคำที่ใช้ในเอกสาร
% เปลี่ยนเป็นภาษาอังกฤษ: แก้ % ---------------------------------------------------------------
% preamble/lang-th.tex - ฟอนต์ไทย การตัดคำ และคำที่ใช้ในเอกสาร
% เปลี่ยนเป็นภาษาอังกฤษ: แก้ % ---------------------------------------------------------------
% preamble/lang-th.tex - ฟอนต์ไทย การตัดคำ และคำที่ใช้ในเอกสาร
% เปลี่ยนเป็นภาษาอังกฤษ: แก้ % ---------------------------------------------------------------
% preamble/lang-th.tex - ฟอนต์ไทย การตัดคำ และคำที่ใช้ในเอกสาร
% เปลี่ยนเป็นภาษาอังกฤษ: แก้ \input{preamble/lang-th} ใน main.tex เป็น lang-en
% ---------------------------------------------------------------
\newif\iflangth
\langthtrue

% \newtxt{\txtfoo}{ไทย}{English}: คำที่ layout ใช้ เลือกภาษาตามไฟล์นี้
\newcommand{\newtxt}[3]{\newcommand{#1}{#2}}

% ฟอนต์หลัก Sarabun (มีอักษรละตินในตัว)
\setmainfont{Sarabun}
\IfFontExistsTF{IBM Plex Sans Thai}{\setsansfont{IBM Plex Sans Thai}}{\setsansfont{Sarabun}}
\IfFontExistsTF{JetBrains Mono}{\setmonofont{JetBrains Mono}[Scale=MatchLowercase, Contextuals=NoAlternate]}{}   % no <= ligatures in code
\setmathfont{latinmodern-math.otf}[Scale=MatchLowercase]
\newfontfamily\thaifallbackfont{Sarabun}   % ตัวไทยในโค้ด/\texttt (ดู macros.tex)

% การตัดบรรทัดภาษาไทย (ไทยไม่มีช่องว่างระหว่างคำ)
\XeTeXlinebreaklocale "th"
\XeTeXlinebreakskip = 0pt plus 0.1em

% ระยะบรรทัด 1.5 เท่า อ่านง่ายสำหรับฟอนต์ไทย
\onehalfspacing

% --- คำในเอกสาร ---
\newcommand{\txttheorem}{ทฤษฎีบท}
\newcommand{\txtlemma}{บทตั้ง}
\newcommand{\txtproposition}{ประพจน์}
\newcommand{\txtcorollary}{บทแทรก}
\newcommand{\txtconjecture}{ข้อความคาดการณ์}
\newcommand{\txtdefinition}{บทนิยาม}
\newcommand{\txtexample}{ตัวอย่าง}
\newcommand{\txtexercise}{แบบฝึกหัด}
\newcommand{\txtremark}{ข้อสังเกต}
\newcommand{\txtnote}{หมายเหตุ}
\newcommand{\txtproof}{พิสูจน์}
\newcommand{\txtsolution}{วิธีทำ}
\newcommand{\txtanswer}{เฉลย}
\newcommand{\txtwork}{แสดงวิธีทำ}
\newcommand{\txtpoints}{คะแนน}
\newcommand{\txtkeypoints}{ข้อควรจำ}
\newcommand{\txtrefs}{เอกสารอ้างอิง}
\newcommand{\txtbib}{บรรณานุกรม}
\newcommand{\txtdraft}{ฉบับร่าง}
\newcommand{\txtst}{\text{ โดยที่ }}
\newcommand{\txtreview}{ทบทวน}
\newcommand{\txtgoals}{สิ่งที่ควรทำได้หลังทบทวน}
\newcommand{\txtsummary}{สรุป}
\newcommand{\txtworked}{ตัวอย่างพร้อมวิธีทำ}
\newcommand{\txtpartmc}{ตอนที่ 1 ปรนัย}
\newcommand{\txtpartwritten}{ตอนที่ 2 อัตนัย}
\newcommand{\txtmistakes}{ข้อผิดพลาดที่พบบ่อย}
\newcommand{\txtselfcheck}{ตรวจสอบตนเอง}
\newcommand{\txtname}{ชื่อ}
\newcommand{\txtnumber}{เลขที่}
\newcommand{\txtstudentid}{รหัสนักศึกษา}
\newcommand{\txtdate}{วันที่}
\newcommand{\txtscore}{คะแนน}
\newcommand{\txtpage}{หน้า}
\newcommand{\txtof}{จาก}
\newcommand{\txtanswerkey}{ฉบับเฉลย}

\renewcommand{\proofname}{\txtproof}
\renewcommand{\contentsname}{สารบัญ}
\renewcommand{\listfigurename}{สารบัญภาพ}
\renewcommand{\listtablename}{สารบัญตาราง}
\renewcommand{\figurename}{ภาพที่}
\renewcommand{\tablename}{ตารางที่}
\def\refname{\txtrefs}
\def\bibname{บรรณานุกรม}
\def\abstractname{บทคัดย่อ}
\def\chaptername{บทที่}
\def\appendixname{ภาคผนวก}

% วันที่แบบไทย พ.ศ.
\renewcommand{\today}{\number\day\space
  \ifcase\month\or มกราคม\or กุมภาพันธ์\or มีนาคม\or เมษายน\or พฤษภาคม\or มิถุนายน\or
  กรกฎาคม\or สิงหาคม\or กันยายน\or ตุลาคม\or พฤศจิกายน\or ธันวาคม\fi
  \space พ.ศ.~\number\numexpr\year+543\relax}

% ตัวนับ ก ข ค ง ... ใช้กับ enumitem ได้เหมือน \alph*
% และใช้เลขหน้าได้: \pagenumbering{thaialph}
\newcommand{\thaialphnum}[1]{%
  \ifcase#1\or ก\or ข\or ค\or ง\or จ\or ฉ\or ช\or ซ\or ฌ\or ญ\or ฎ\or ฏ\or
  ฐ\or ฑ\or ฒ\or ณ\or ด\or ต\or ถ\or ท\or ธ\or น\or บ\or ป\or ผ\or ฝ\or
  พ\or ฟ\or ภ\or ม\or ย\or ร\or ล\or ว\or ศ\or ษ\or ส\or ห\or ฬ\or อ\or ฮ\fi}
\makeatletter
\newcommand{\thaialph}[1]{\expandafter\thaialphnum\csname c@#1\endcsname}
\def\@thaialph#1{\thaialphnum{#1}}
\makeatother
\AddEnumerateCounter{\thaialph}{\thaialphnum}{ญ}

% --- ชื่อที่ \cref ใช้ ---
\crefname{equation}{สมการ}{สมการ}
\creflabelformat{equation}{#2(#1)#3}
\crefname{chapter}{บทที่}{บทที่}
\crefname{appendix}{ภาคผนวก}{ภาคผนวก}
\crefname{section}{หัวข้อ}{หัวข้อ}
\crefname{subsection}{หัวข้อ}{หัวข้อ}
\crefname{figure}{ภาพที่}{ภาพที่}
\crefname{table}{ตารางที่}{ตารางที่}
\crefname{theorem}{\txttheorem}{\txttheorem}
\crefname{lemma}{\txtlemma}{\txtlemma}
\crefname{proposition}{\txtproposition}{\txtproposition}
\crefname{corollary}{\txtcorollary}{\txtcorollary}
\crefname{conjecture}{\txtconjecture}{\txtconjecture}
\crefname{definition}{\txtdefinition}{\txtdefinition}
\crefname{example}{\txtexample}{\txtexample}
\crefname{exercise}{\txtexercise}{\txtexercise}
\crefname{remark}{\txtremark}{\txtremark}
\crefname{enumi}{ข้อ}{ข้อ}
\crefname{exercisesi}{ข้อ}{ข้อ}
\newcommand{\crefpairconjunction}{ และ }
\newcommand{\crefmiddleconjunction}{, }
\newcommand{\creflastconjunction}{ และ }
\newcommand{\crefrangeconjunction}{ ถึง }
 ใน main.tex เป็น lang-en
% ---------------------------------------------------------------
\newif\iflangth
\langthtrue

% \newtxt{\txtfoo}{ไทย}{English}: คำที่ layout ใช้ เลือกภาษาตามไฟล์นี้
\newcommand{\newtxt}[3]{\newcommand{#1}{#2}}

% ฟอนต์หลัก Sarabun (มีอักษรละตินในตัว)
\setmainfont{Sarabun}
\IfFontExistsTF{IBM Plex Sans Thai}{\setsansfont{IBM Plex Sans Thai}}{\setsansfont{Sarabun}}
\IfFontExistsTF{JetBrains Mono}{\setmonofont{JetBrains Mono}[Scale=MatchLowercase, Contextuals=NoAlternate]}{}   % no <= ligatures in code
\setmathfont{latinmodern-math.otf}[Scale=MatchLowercase]
\newfontfamily\thaifallbackfont{Sarabun}   % ตัวไทยในโค้ด/\texttt (ดู macros.tex)

% การตัดบรรทัดภาษาไทย (ไทยไม่มีช่องว่างระหว่างคำ)
\XeTeXlinebreaklocale "th"
\XeTeXlinebreakskip = 0pt plus 0.1em

% ระยะบรรทัด 1.5 เท่า อ่านง่ายสำหรับฟอนต์ไทย
\onehalfspacing

% --- คำในเอกสาร ---
\newcommand{\txttheorem}{ทฤษฎีบท}
\newcommand{\txtlemma}{บทตั้ง}
\newcommand{\txtproposition}{ประพจน์}
\newcommand{\txtcorollary}{บทแทรก}
\newcommand{\txtconjecture}{ข้อความคาดการณ์}
\newcommand{\txtdefinition}{บทนิยาม}
\newcommand{\txtexample}{ตัวอย่าง}
\newcommand{\txtexercise}{แบบฝึกหัด}
\newcommand{\txtremark}{ข้อสังเกต}
\newcommand{\txtnote}{หมายเหตุ}
\newcommand{\txtproof}{พิสูจน์}
\newcommand{\txtsolution}{วิธีทำ}
\newcommand{\txtanswer}{เฉลย}
\newcommand{\txtwork}{แสดงวิธีทำ}
\newcommand{\txtpoints}{คะแนน}
\newcommand{\txtkeypoints}{ข้อควรจำ}
\newcommand{\txtrefs}{เอกสารอ้างอิง}
\newcommand{\txtbib}{บรรณานุกรม}
\newcommand{\txtdraft}{ฉบับร่าง}
\newcommand{\txtst}{\text{ โดยที่ }}
\newcommand{\txtreview}{ทบทวน}
\newcommand{\txtgoals}{สิ่งที่ควรทำได้หลังทบทวน}
\newcommand{\txtsummary}{สรุป}
\newcommand{\txtworked}{ตัวอย่างพร้อมวิธีทำ}
\newcommand{\txtpartmc}{ตอนที่ 1 ปรนัย}
\newcommand{\txtpartwritten}{ตอนที่ 2 อัตนัย}
\newcommand{\txtmistakes}{ข้อผิดพลาดที่พบบ่อย}
\newcommand{\txtselfcheck}{ตรวจสอบตนเอง}
\newcommand{\txtname}{ชื่อ}
\newcommand{\txtnumber}{เลขที่}
\newcommand{\txtstudentid}{รหัสนักศึกษา}
\newcommand{\txtdate}{วันที่}
\newcommand{\txtscore}{คะแนน}
\newcommand{\txtpage}{หน้า}
\newcommand{\txtof}{จาก}
\newcommand{\txtanswerkey}{ฉบับเฉลย}

\renewcommand{\proofname}{\txtproof}
\renewcommand{\contentsname}{สารบัญ}
\renewcommand{\listfigurename}{สารบัญภาพ}
\renewcommand{\listtablename}{สารบัญตาราง}
\renewcommand{\figurename}{ภาพที่}
\renewcommand{\tablename}{ตารางที่}
\def\refname{\txtrefs}
\def\bibname{บรรณานุกรม}
\def\abstractname{บทคัดย่อ}
\def\chaptername{บทที่}
\def\appendixname{ภาคผนวก}

% วันที่แบบไทย พ.ศ.
\renewcommand{\today}{\number\day\space
  \ifcase\month\or มกราคม\or กุมภาพันธ์\or มีนาคม\or เมษายน\or พฤษภาคม\or มิถุนายน\or
  กรกฎาคม\or สิงหาคม\or กันยายน\or ตุลาคม\or พฤศจิกายน\or ธันวาคม\fi
  \space พ.ศ.~\number\numexpr\year+543\relax}

% ตัวนับ ก ข ค ง ... ใช้กับ enumitem ได้เหมือน \alph*
% และใช้เลขหน้าได้: \pagenumbering{thaialph}
\newcommand{\thaialphnum}[1]{%
  \ifcase#1\or ก\or ข\or ค\or ง\or จ\or ฉ\or ช\or ซ\or ฌ\or ญ\or ฎ\or ฏ\or
  ฐ\or ฑ\or ฒ\or ณ\or ด\or ต\or ถ\or ท\or ธ\or น\or บ\or ป\or ผ\or ฝ\or
  พ\or ฟ\or ภ\or ม\or ย\or ร\or ล\or ว\or ศ\or ษ\or ส\or ห\or ฬ\or อ\or ฮ\fi}
\makeatletter
\newcommand{\thaialph}[1]{\expandafter\thaialphnum\csname c@#1\endcsname}
\def\@thaialph#1{\thaialphnum{#1}}
\makeatother
\AddEnumerateCounter{\thaialph}{\thaialphnum}{ญ}

% --- ชื่อที่ \cref ใช้ ---
\crefname{equation}{สมการ}{สมการ}
\creflabelformat{equation}{#2(#1)#3}
\crefname{chapter}{บทที่}{บทที่}
\crefname{appendix}{ภาคผนวก}{ภาคผนวก}
\crefname{section}{หัวข้อ}{หัวข้อ}
\crefname{subsection}{หัวข้อ}{หัวข้อ}
\crefname{figure}{ภาพที่}{ภาพที่}
\crefname{table}{ตารางที่}{ตารางที่}
\crefname{theorem}{\txttheorem}{\txttheorem}
\crefname{lemma}{\txtlemma}{\txtlemma}
\crefname{proposition}{\txtproposition}{\txtproposition}
\crefname{corollary}{\txtcorollary}{\txtcorollary}
\crefname{conjecture}{\txtconjecture}{\txtconjecture}
\crefname{definition}{\txtdefinition}{\txtdefinition}
\crefname{example}{\txtexample}{\txtexample}
\crefname{exercise}{\txtexercise}{\txtexercise}
\crefname{remark}{\txtremark}{\txtremark}
\crefname{enumi}{ข้อ}{ข้อ}
\crefname{exercisesi}{ข้อ}{ข้อ}
\newcommand{\crefpairconjunction}{ และ }
\newcommand{\crefmiddleconjunction}{, }
\newcommand{\creflastconjunction}{ และ }
\newcommand{\crefrangeconjunction}{ ถึง }
 ใน main.tex เป็น lang-en
% ---------------------------------------------------------------
\newif\iflangth
\langthtrue

% \newtxt{\txtfoo}{ไทย}{English}: คำที่ layout ใช้ เลือกภาษาตามไฟล์นี้
\newcommand{\newtxt}[3]{\newcommand{#1}{#2}}

% ฟอนต์หลัก Sarabun (มีอักษรละตินในตัว)
\setmainfont{Sarabun}
\IfFontExistsTF{IBM Plex Sans Thai}{\setsansfont{IBM Plex Sans Thai}}{\setsansfont{Sarabun}}
\IfFontExistsTF{JetBrains Mono}{\setmonofont{JetBrains Mono}[Scale=MatchLowercase, Contextuals=NoAlternate]}{}   % no <= ligatures in code
\setmathfont{latinmodern-math.otf}[Scale=MatchLowercase]
\newfontfamily\thaifallbackfont{Sarabun}   % ตัวไทยในโค้ด/\texttt (ดู macros.tex)

% การตัดบรรทัดภาษาไทย (ไทยไม่มีช่องว่างระหว่างคำ)
\XeTeXlinebreaklocale "th"
\XeTeXlinebreakskip = 0pt plus 0.1em

% ระยะบรรทัด 1.5 เท่า อ่านง่ายสำหรับฟอนต์ไทย
\onehalfspacing

% --- คำในเอกสาร ---
\newcommand{\txttheorem}{ทฤษฎีบท}
\newcommand{\txtlemma}{บทตั้ง}
\newcommand{\txtproposition}{ประพจน์}
\newcommand{\txtcorollary}{บทแทรก}
\newcommand{\txtconjecture}{ข้อความคาดการณ์}
\newcommand{\txtdefinition}{บทนิยาม}
\newcommand{\txtexample}{ตัวอย่าง}
\newcommand{\txtexercise}{แบบฝึกหัด}
\newcommand{\txtremark}{ข้อสังเกต}
\newcommand{\txtnote}{หมายเหตุ}
\newcommand{\txtproof}{พิสูจน์}
\newcommand{\txtsolution}{วิธีทำ}
\newcommand{\txtanswer}{เฉลย}
\newcommand{\txtwork}{แสดงวิธีทำ}
\newcommand{\txtpoints}{คะแนน}
\newcommand{\txtkeypoints}{ข้อควรจำ}
\newcommand{\txtrefs}{เอกสารอ้างอิง}
\newcommand{\txtbib}{บรรณานุกรม}
\newcommand{\txtdraft}{ฉบับร่าง}
\newcommand{\txtst}{\text{ โดยที่ }}
\newcommand{\txtreview}{ทบทวน}
\newcommand{\txtgoals}{สิ่งที่ควรทำได้หลังทบทวน}
\newcommand{\txtsummary}{สรุป}
\newcommand{\txtworked}{ตัวอย่างพร้อมวิธีทำ}
\newcommand{\txtpartmc}{ตอนที่ 1 ปรนัย}
\newcommand{\txtpartwritten}{ตอนที่ 2 อัตนัย}
\newcommand{\txtmistakes}{ข้อผิดพลาดที่พบบ่อย}
\newcommand{\txtselfcheck}{ตรวจสอบตนเอง}
\newcommand{\txtname}{ชื่อ}
\newcommand{\txtnumber}{เลขที่}
\newcommand{\txtstudentid}{รหัสนักศึกษา}
\newcommand{\txtdate}{วันที่}
\newcommand{\txtscore}{คะแนน}
\newcommand{\txtpage}{หน้า}
\newcommand{\txtof}{จาก}
\newcommand{\txtanswerkey}{ฉบับเฉลย}

\renewcommand{\proofname}{\txtproof}
\renewcommand{\contentsname}{สารบัญ}
\renewcommand{\listfigurename}{สารบัญภาพ}
\renewcommand{\listtablename}{สารบัญตาราง}
\renewcommand{\figurename}{ภาพที่}
\renewcommand{\tablename}{ตารางที่}
\def\refname{\txtrefs}
\def\bibname{บรรณานุกรม}
\def\abstractname{บทคัดย่อ}
\def\chaptername{บทที่}
\def\appendixname{ภาคผนวก}

% วันที่แบบไทย พ.ศ.
\renewcommand{\today}{\number\day\space
  \ifcase\month\or มกราคม\or กุมภาพันธ์\or มีนาคม\or เมษายน\or พฤษภาคม\or มิถุนายน\or
  กรกฎาคม\or สิงหาคม\or กันยายน\or ตุลาคม\or พฤศจิกายน\or ธันวาคม\fi
  \space พ.ศ.~\number\numexpr\year+543\relax}

% ตัวนับ ก ข ค ง ... ใช้กับ enumitem ได้เหมือน \alph*
% และใช้เลขหน้าได้: \pagenumbering{thaialph}
\newcommand{\thaialphnum}[1]{%
  \ifcase#1\or ก\or ข\or ค\or ง\or จ\or ฉ\or ช\or ซ\or ฌ\or ญ\or ฎ\or ฏ\or
  ฐ\or ฑ\or ฒ\or ณ\or ด\or ต\or ถ\or ท\or ธ\or น\or บ\or ป\or ผ\or ฝ\or
  พ\or ฟ\or ภ\or ม\or ย\or ร\or ล\or ว\or ศ\or ษ\or ส\or ห\or ฬ\or อ\or ฮ\fi}
\makeatletter
\newcommand{\thaialph}[1]{\expandafter\thaialphnum\csname c@#1\endcsname}
\def\@thaialph#1{\thaialphnum{#1}}
\makeatother
\AddEnumerateCounter{\thaialph}{\thaialphnum}{ญ}

% --- ชื่อที่ \cref ใช้ ---
\crefname{equation}{สมการ}{สมการ}
\creflabelformat{equation}{#2(#1)#3}
\crefname{chapter}{บทที่}{บทที่}
\crefname{appendix}{ภาคผนวก}{ภาคผนวก}
\crefname{section}{หัวข้อ}{หัวข้อ}
\crefname{subsection}{หัวข้อ}{หัวข้อ}
\crefname{figure}{ภาพที่}{ภาพที่}
\crefname{table}{ตารางที่}{ตารางที่}
\crefname{theorem}{\txttheorem}{\txttheorem}
\crefname{lemma}{\txtlemma}{\txtlemma}
\crefname{proposition}{\txtproposition}{\txtproposition}
\crefname{corollary}{\txtcorollary}{\txtcorollary}
\crefname{conjecture}{\txtconjecture}{\txtconjecture}
\crefname{definition}{\txtdefinition}{\txtdefinition}
\crefname{example}{\txtexample}{\txtexample}
\crefname{exercise}{\txtexercise}{\txtexercise}
\crefname{remark}{\txtremark}{\txtremark}
\crefname{enumi}{ข้อ}{ข้อ}
\crefname{exercisesi}{ข้อ}{ข้อ}
\newcommand{\crefpairconjunction}{ และ }
\newcommand{\crefmiddleconjunction}{, }
\newcommand{\creflastconjunction}{ และ }
\newcommand{\crefrangeconjunction}{ ถึง }
 ใน main.tex เป็น lang-en
% ---------------------------------------------------------------
\newif\iflangth
\langthtrue

% \newtxt{\txtfoo}{ไทย}{English}: คำที่ layout ใช้ เลือกภาษาตามไฟล์นี้
\newcommand{\newtxt}[3]{\newcommand{#1}{#2}}

% ฟอนต์หลัก Sarabun (มีอักษรละตินในตัว)
\setmainfont{Sarabun}
\IfFontExistsTF{IBM Plex Sans Thai}{\setsansfont{IBM Plex Sans Thai}}{\setsansfont{Sarabun}}
\IfFontExistsTF{JetBrains Mono}{\setmonofont{JetBrains Mono}[Scale=MatchLowercase, Contextuals=NoAlternate]}{}   % no <= ligatures in code
\setmathfont{latinmodern-math.otf}[Scale=MatchLowercase]
\newfontfamily\thaifallbackfont{Sarabun}   % ตัวไทยในโค้ด/\texttt (ดู macros.tex)

% การตัดบรรทัดภาษาไทย (ไทยไม่มีช่องว่างระหว่างคำ)
\XeTeXlinebreaklocale "th"
\XeTeXlinebreakskip = 0pt plus 0.1em

% ระยะบรรทัด 1.5 เท่า อ่านง่ายสำหรับฟอนต์ไทย
\onehalfspacing

% --- คำในเอกสาร ---
\newcommand{\txttheorem}{ทฤษฎีบท}
\newcommand{\txtlemma}{บทตั้ง}
\newcommand{\txtproposition}{ประพจน์}
\newcommand{\txtcorollary}{บทแทรก}
\newcommand{\txtconjecture}{ข้อความคาดการณ์}
\newcommand{\txtdefinition}{บทนิยาม}
\newcommand{\txtexample}{ตัวอย่าง}
\newcommand{\txtexercise}{แบบฝึกหัด}
\newcommand{\txtremark}{ข้อสังเกต}
\newcommand{\txtnote}{หมายเหตุ}
\newcommand{\txtproof}{พิสูจน์}
\newcommand{\txtsolution}{วิธีทำ}
\newcommand{\txtanswer}{เฉลย}
\newcommand{\txtwork}{แสดงวิธีทำ}
\newcommand{\txtpoints}{คะแนน}
\newcommand{\txtkeypoints}{ข้อควรจำ}
\newcommand{\txtrefs}{เอกสารอ้างอิง}
\newcommand{\txtbib}{บรรณานุกรม}
\newcommand{\txtdraft}{ฉบับร่าง}
\newcommand{\txtst}{\text{ โดยที่ }}
\newcommand{\txtreview}{ทบทวน}
\newcommand{\txtgoals}{สิ่งที่ควรทำได้หลังทบทวน}
\newcommand{\txtsummary}{สรุป}
\newcommand{\txtworked}{ตัวอย่างพร้อมวิธีทำ}
\newcommand{\txtpartmc}{ตอนที่ 1 ปรนัย}
\newcommand{\txtpartwritten}{ตอนที่ 2 อัตนัย}
\newcommand{\txtmistakes}{ข้อผิดพลาดที่พบบ่อย}
\newcommand{\txtselfcheck}{ตรวจสอบตนเอง}
\newcommand{\txtname}{ชื่อ}
\newcommand{\txtnumber}{เลขที่}
\newcommand{\txtstudentid}{รหัสนักศึกษา}
\newcommand{\txtdate}{วันที่}
\newcommand{\txtscore}{คะแนน}
\newcommand{\txtpage}{หน้า}
\newcommand{\txtof}{จาก}
\newcommand{\txtanswerkey}{ฉบับเฉลย}

\renewcommand{\proofname}{\txtproof}
\renewcommand{\contentsname}{สารบัญ}
\renewcommand{\listfigurename}{สารบัญภาพ}
\renewcommand{\listtablename}{สารบัญตาราง}
\renewcommand{\figurename}{ภาพที่}
\renewcommand{\tablename}{ตารางที่}
\def\refname{\txtrefs}
\def\bibname{บรรณานุกรม}
\def\abstractname{บทคัดย่อ}
\def\chaptername{บทที่}
\def\appendixname{ภาคผนวก}

% วันที่แบบไทย พ.ศ.
\renewcommand{\today}{\number\day\space
  \ifcase\month\or มกราคม\or กุมภาพันธ์\or มีนาคม\or เมษายน\or พฤษภาคม\or มิถุนายน\or
  กรกฎาคม\or สิงหาคม\or กันยายน\or ตุลาคม\or พฤศจิกายน\or ธันวาคม\fi
  \space พ.ศ.~\number\numexpr\year+543\relax}

% ตัวนับ ก ข ค ง ... ใช้กับ enumitem ได้เหมือน \alph*
% และใช้เลขหน้าได้: \pagenumbering{thaialph}
\newcommand{\thaialphnum}[1]{%
  \ifcase#1\or ก\or ข\or ค\or ง\or จ\or ฉ\or ช\or ซ\or ฌ\or ญ\or ฎ\or ฏ\or
  ฐ\or ฑ\or ฒ\or ณ\or ด\or ต\or ถ\or ท\or ธ\or น\or บ\or ป\or ผ\or ฝ\or
  พ\or ฟ\or ภ\or ม\or ย\or ร\or ล\or ว\or ศ\or ษ\or ส\or ห\or ฬ\or อ\or ฮ\fi}
\makeatletter
\newcommand{\thaialph}[1]{\expandafter\thaialphnum\csname c@#1\endcsname}
\def\@thaialph#1{\thaialphnum{#1}}
\makeatother
\AddEnumerateCounter{\thaialph}{\thaialphnum}{ญ}

% --- ชื่อที่ \cref ใช้ ---
\crefname{equation}{สมการ}{สมการ}
\creflabelformat{equation}{#2(#1)#3}
\crefname{chapter}{บทที่}{บทที่}
\crefname{appendix}{ภาคผนวก}{ภาคผนวก}
\crefname{section}{หัวข้อ}{หัวข้อ}
\crefname{subsection}{หัวข้อ}{หัวข้อ}
\crefname{figure}{ภาพที่}{ภาพที่}
\crefname{table}{ตารางที่}{ตารางที่}
\crefname{theorem}{\txttheorem}{\txttheorem}
\crefname{lemma}{\txtlemma}{\txtlemma}
\crefname{proposition}{\txtproposition}{\txtproposition}
\crefname{corollary}{\txtcorollary}{\txtcorollary}
\crefname{conjecture}{\txtconjecture}{\txtconjecture}
\crefname{definition}{\txtdefinition}{\txtdefinition}
\crefname{example}{\txtexample}{\txtexample}
\crefname{exercise}{\txtexercise}{\txtexercise}
\crefname{remark}{\txtremark}{\txtremark}
\crefname{enumi}{ข้อ}{ข้อ}
\crefname{exercisesi}{ข้อ}{ข้อ}
\newcommand{\crefpairconjunction}{ และ }
\newcommand{\crefmiddleconjunction}{, }
\newcommand{\creflastconjunction}{ และ }
\newcommand{\crefrangeconjunction}{ ถึง }
